\documentclass[11pt]{article}
\usepackage[a4paper,margin=1in]{geometry}
\usepackage{graphicx}
\usepackage{booktabs}
\usepackage{amsmath}
\usepackage{amssymb}
\usepackage{hyperref}
\usepackage{natbib}
\usepackage{xcolor}
\usepackage{caption}
\usepackage{float}

\title{Digital telomere measurement by long-read sequencing distinguishes healthy aging from disease}
\author{}
\date{}

\begin{document}
\maketitle

\begin{abstract}
Telomere length is an established biomarker of organismal aging and cellular
replicative potential, yet existing assays---TRF Southern blot, flow-FISH,
qPCR, STELA and TeSLA---return either a coarse mean length, a relative
telomere content, or a view biased toward the shortest fraction. None
delivers the full single-molecule length distribution at the throughput
and precision required for clinical use. Here we present digital telomere
measurement by nanopore long-read sequencing, combining a telomere-capture
oligonucleotide library preparation with the Telometer pipeline, which
anchors each intact telomeric read to a specific chromosome end on the
telomere-to-telomere (T2T) human reference. The method reaches a plateau
precision of 30--40 bp, correlates with TRF ($n{=}14$) and flow-FISH
($n{=}7$, $R^{2}{=}0.84$ in RTEL1-mutant samples), detects $\sim$40 bp/day
of attrition in TERT-null hESCs, and measures $\sim$1000 bp of de novo
elongation in 3 days of telomerase overexpression. In peripheral blood
leukocytes from 14 healthy donors aged 18--77, mean and median lengths
decline by $\sim$24 bp/year, and the third quartile shortens faster than
the first. A classifier built on the telomere distribution and age
distinguishes healthy donors from symptomatic patients (AUC $=$ 0.98),
asymptomatic telomere biology disorder carriers (AUC $=$ 0.93) and the
combined group (AUC $=$ 0.91). In 20 colorectal carcinoma patients, tumour
telomeres are shorter than matched benign epithelia in 75\% of cases.
Digital telomere measurement turns the full length distribution into a
clinical biomarker platform for aging and disease.
\end{abstract}

\section{Introduction}

Telomeres are the nucleoprotein caps that guarantee the integrity of
linear eukaryotic chromosomes, solving the end-replication problem and
preventing chromosome termini from being misread as double-strand breaks
\citep{delange2005shelterin,palm2008howshelterin}. In human somatic cells,
the incremental loss of terminal TTAGGG repeats at every division
eventually triggers replicative senescence, and this attritional clock
underlies much of what we recognise as cellular and organismal aging
\citep{harley1990telomeres,blackburn2006human}. Because telomere length
integrates replicative history with the activity of telomerase and the
shelterin complex, it has long been used as a biomarker of aging, of
risk for age-associated disease, and of rare inherited disorders of
telomere maintenance collectively known as telomere biology disorders
(TBDs) \citep{armanios2012telomeres,bertuch2016molecular,aubert2012telomere}.

The practical value of this biomarker, however, is held back by the
methods used to measure it. Terminal restriction fragment (TRF)
Southern blotting reports a bulk mean length as a smear on a gel; qPCR
and flow-FISH estimate relative telomeric signal rather than the length
of any individual telomere
\citep{aubert2012telomere,baerlocher2006flow,cawthon2002telomere}.
STELA and TeSLA resolve individual chromosome ends using PCR
amplification but are bounded by polymerase processivity and tend to
over-represent the shortest fraction of the distribution
\citep{baird2003extensive,lai2018comparison}. Short-read sequencing
methods such as TelSeq extract a relative telomeric content from
whole-genome data and have been applied to cohorts such as TCGA, but
they cannot recover the length of any single telomere
\citep{ding2014telseq,barthel2017systematic,demanelis2020determinants}.
Long-read platforms such as Oxford Nanopore and PacBio HiFi, together
with the recently completed telomere-to-telomere (T2T) human genome, now
make it possible in principle to measure individual intact telomeres
with chromosome-level anchoring
\citep{nurk2022complete,tham2023comprehensive,schmidt2023targeted,karimian2024human,heller2019telomeres},
but no pipeline has yet united these advances into a clinical workflow
capable of distinguishing healthy aging from disease from the full
length distribution.

Here we close this gap with a digital telomere measurement pipeline
built around nanopore long-read sequencing. A biotinylated
oligonucleotide complementary to the 3$'$ overhang, combined with
restriction digestion of non-telomeric genomic DNA, enriches telomeric
reads by several thousand-fold while leaving the length distribution
undistorted. A companion bioinformatic tool, Telometer, aligns telomeric
repeats to chromosomal termini of the T2T reference and defines the
extent of each telomere as the span between the first terminal TTAGGG
repeat and the final two consecutive repeats before the subtelomere.
The output is a chromosome-anchored length for every intact telomeric
read in the sample, from which per-chromosome distributions, summary
statistics and classifiers can be derived. We validate the method against
TRF ($n{=}14$) and flow-FISH ($n{=}7$), quantify its precision by
bootstrapping, and apply it to genetically defined hESC models of
telomere dysfunction, to peripheral blood leukocytes from 14 healthy
donors spanning ages 18--77, to RTEL1-mutant individuals and a Stanford
patient with a newly discovered TINF2 frameshift mutation, and to
patient-matched benign and malignant colon tissue from 20 colorectal
carcinoma cases
\citep{walne2013constitutional,tummala2015poly,savage2008tinf2,alder2018diagnostic}.

Our main contributions are as follows. First, we establish that digital
telomere measurement plateaus at a precision of 30--40 bp, correlates
tightly with TRF and flow-FISH, and show that TRF systematically
overestimates mean telomere length by 1000--3000 bp. Second, we show
that the pipeline recovers $\sim$40 bp/day of attrition in TERT-null
hESCs and $\sim$1000 bp of de novo elongation after three days of
telomerase overexpression, with chromosome-anchored evidence that
telomerase preferentially extends the chromosomes with shortest baseline
telomeres
\citep{hemann2001shortest,britt2012preferential,teixeira2004intertelomere}.
Third, in peripheral blood leukocytes from 14 healthy donors, mean and
median telomere length decline by $\sim$24 bp/year, and the third
quartile shortens faster than the first---an asymmetry abolished in
RTEL1 carriers and in TERT-null hESCs. Fourth, a simple machine-learning
classifier trained on the length distribution and donor age distinguishes
healthy donors from symptomatic TBD patients (AUC $=$ 0.98), from
asymptomatic carriers that flow-FISH cannot detect (AUC $=$ 0.93), and
from the combined group (AUC $=$ 0.91). Finally, in matched benign and
malignant colorectal tissue, tumour telomeres are shorter than benign
epithelium in 75\% of the 20 patients and show no correlation with donor
age, consistent with telomerase-stabilised tumour distributions
\citep{barthel2017systematic,roka2023tumor}.

\section{Related Work}

Analog telomere measurement has, for decades, relied on TRF, flow-FISH,
Q-FISH and qPCR, each returning a coarse mean or relative signal rather
than individual telomere lengths
\citep{aubert2012telomere,baerlocher2006flow,cawthon2002telomere,lai2017comparison}.
TRF is further limited by incomplete digestion of subtelomeric DNA,
which inflates the apparent mean, while qPCR reports only a ratio to a
single-copy gene. Flow-FISH is the clinical standard for peripheral
blood leukocytes (PBLs) and a central element of the diagnostic pathway
for TBDs, but it cannot distinguish asymptomatic TBD carriers from
healthy donors in a consistent way \citep{alder2018diagnostic}.

PCR-based single-telomere assays sharpen the picture. STELA ligates a
linker to the G-overhang and amplifies a tagged terminal fragment, first
revealing extensive allelic variation and ultrashort telomeres at
senescence \citep{baird2003extensive}. TeSLA extends this strategy
genome-wide by restriction digestion, adapter ligation and amplification
of short telomeric fragments, making the assay sensitive to the shortest
fraction \citep{lai2018comparison}. Both methods, however, are capped by
polymerase processivity and systematically under-report long telomeres.

Sequencing-based approaches bridge this limitation partially. Short-read
tools such as TelSeq count reads containing telomeric repeats in
whole-genome data and have been used at cohort scale
\citep{ding2014telseq,barthel2017systematic,demanelis2020determinants}.
They cannot recover the length of an individual telomere and cannot
separate changes in distribution structure from changes in total
telomeric content. Long-read sequencing technologies, paired with the
T2T reference, have begun to offer single-molecule, per-chromosome
measurements
\citep{tham2023comprehensive,schmidt2023targeted,karimian2024human,nurk2022complete,heller2019telomeres},
but published pipelines have either been limited to whole-genome
sequencing, which is costly per sample, or to highly specialised
platforms, without joining these advances to aging and disease
diagnostics.

Telomere biology disorders provide both the motivation and the benchmark
for our work. Dyskeratosis congenita and related syndromes are caused by
germline defects in telomerase components (TERT, TERC, PARN, DKC1), in
shelterin (TINF2) or in the helicase RTEL1
\citep{tummala2015poly,walne2013constitutional,savage2008tinf2,armanios2012telomeres,bertuch2016molecular}.
Clinical triage combines flow-FISH, population norms and targeted exon
sequencing \citep{alder2018diagnostic}. Because TBDs display genetic
anticipation and because asymptomatic carriers often fall in the lower
quintile of the age-matched telomere distribution, a higher-resolution
assay capable of detecting carrier status is a clear unmet need.

Finally, length regulation itself is a conserved, measurable biology.
Telomerase preferentially elongates the shortest telomeres in yeast and
in mammalian systems
\citep{teixeira2004intertelomere,hemann2001shortest,britt2012preferential},
and shelterin components modulate telomerase access to act as a
negative-feedback rheostat on length
\citep{palm2008howshelterin,hockemeyer2015control}. In cancer,
telomerase reactivation stabilises tumour length distributions at short
but persistent set-points \citep{chiba2015cancer,barthel2017systematic}.
Our single-molecule measurements make these rheostat-level regulatory
signatures directly readable from the interquartile structure of the
distribution.

\section{Methods}

\subsection{Samples and cohorts}

Five biological sources were sequenced. (i) HEK293T cells used for
head-to-head comparison of telomere capture versus whole-genome
sequencing and for transient-transfection de novo elongation assays.
(ii) Genetically modified human embryonic stem cells (hESCs), including
wild-type, PARN knock-out, TIN2 T284R heterozygous and homozygous
mutants, and TERT-inactivated cells (CDKN2A$^{-/-}$ TERT$^{-/-}$,
AAVSI:hTERT$^{\text{flox}}$) passaged for 66, 78, 98 and 105 days after
Cre-mediated inactivation. (iii) Peripheral blood leukocytes (PBLs) from
14 healthy human donors aged 18--77 years. (iv) PBLs from 6 individuals
carrying RTEL1 mutations (7 samples, ages 6--50) together with paired
peripheral blood (TB32-B) and bone marrow biopsy (TB32-M) samples from a
Stanford patient with chronic myelomonocytic leukaemia (CMML) and
pulmonary fibrosis who harbours a previously unreported TINF2
frameshift, p.Gln298fs. (v) Previously published 30X whole-genome
long-read sequencing data from 20 patients with matched benign colonic
epithelium and colorectal carcinoma, aged predominantly 50--70.

\subsection{Telomere capture library preparation}

Telomere capture uses a biotinylated oligonucleotide complementary to the
3$'$ single-strand overhang. For the canonical human repeat we use
5$'$-TTAGGG-3$'$; for experiments with the TSQ1 variant template we use
5$'$-TTGCGG-3$'$. The oligonucleotide is designed such that one end
hybridises to the telomeric overhang and the other contains an adapter
compatible with the nanopore sequencing library. Genomic DNA is first
restriction-digested to reduce non-telomeric content, the capture
oligonucleotide is annealed, and the telomere-enriched fragments are
prepared as an Oxford Nanopore sequencing library. Variant-specific
capture from genomic DNA of cells not expressing TSQ1 did not yield a
viable library, consistent with the absence of variant repeats in these
samples.

\subsection{The Telometer bioinformatic pipeline}

Telometer identifies telomeric reads by aligning telomeric repeats to
the chromosomal termini of the T2T human reference genome
\citep{nurk2022complete}. The extent of each telomere is defined as the
span from the first terminal repeat at the chromosome terminus to the
last two consecutive TTAGGG repeats before the subtelomeric anchor
sequence. This definition enables chromosome-anchored per-read lengths
from either whole-genome or telomere-enriched libraries. Because
whole-genome sequencing yields few telomeric reads per gigabase, the
capture workflow increases telomere yield by several thousand-fold
without distorting the per-read length distribution.

\subsection{Cell culture and de novo elongation}

TERT knock-out hESCs (CDKN2A$^{-/-}$ TERT$^{-/-}$ AAVSI:hTERT$^{\text{flox}}$)
were passaged for 66, 78, 98 and 105 days after Cre-mediated TERT
inactivation to measure attrition in telomerase-null conditions. De novo
telomere addition was assayed by transiently transfecting HEK293T cells
with plasmids expressing wild-type human telomerase RNA (hTR), the TSQ1
variant template, or GFP, each together with TERT. Telomerase activity
was quantified by the TRAP assay. Three days after transfection,
genomic DNA was harvested and processed by telomere capture sequencing
using either canonical or TSQ1-variant capture oligonucleotides.

\subsection{Reference assays}

TRF Southern blots were performed on 12 of the 14 healthy donors, on
the two TB32 samples, and on genetically modified hESCs. Flow-FISH
was performed on 7 RTEL1-mutant samples. TRF and flow-FISH values were
compared to the digital mean telomere length measured from the
nanopore pipeline.

\subsection{Statistical analyses}

Between-group telomere length comparisons use the Wilcoxon rank-sum
test; this is applied for the PARN-KO vs wild-type contrast in hESCs
($p < 2 \times 10^{-16}$), for aged-cohort contrasts in PBLs (young
18--20, $n{=}5$; middle 35--65, $n{=}6$; elder $>$70, $n{=}3$), and for
paired benign-vs-tumour contrasts in the colorectal carcinoma cohort.
Linear regression is used to relate summary statistics (mean, median,
Q1, Q3) of the telomere length distribution to donor age or days post
TERT inactivation. Precision is estimated by bootstrapping the standard
error of the mean telomere length as a function of the number of
telomeres measured. Method-to-method agreement is reported as the
squared Pearson correlation coefficient ($R^{2}$).

\subsection{Machine-learning classifier}

Telomere measurements and donor ages were randomly and iteratively
assigned to training or test partitions at a 7:3 ratio. Features
comprise per-sample summary statistics of the telomere length
distribution together with donor age. Three binary classifiers were
trained: (i) healthy donors versus symptomatic TBD patients; (ii)
healthy donors versus asymptomatic TBD carriers; and (iii) healthy
donors versus the combined group of symptomatic patients and
asymptomatic carriers. Performance is reported as accuracy on the
held-out test partition and as the area under the receiver operating
characteristic curve (AUC). These three comparisons reflect three
practical diagnostic scenarios in TBD triage: detecting a symptomatic
patient whose phenotype is already clinically apparent, detecting an
asymptomatic carrier that current flow-FISH-based workups frequently
miss, and triaging the full population of individuals at risk for TBD.

\section{Results}

\subsection{Digital telomere measurement by long-read sequencing measures intact telomeres at high resolution}

Existing telomere assays either return a bulk mean or are biased toward
the shortest fragments \citep{aubert2012telomere,baird2003extensive,lai2018comparison}.
To recover the full single-molecule length distribution, we built a
nanopore sequencing and bioinformatic pipeline (Figure~\ref{fig:schematic}).
Reads were aligned with Telometer to the chromosomal termini of the T2T
reference \citep{nurk2022complete}, and each telomere was defined from
its first terminal TTAGGG repeat to the last two consecutive repeats
before the subtelomere, producing a chromosome-anchored per-read length
for every intact telomeric read.

The whole-genome library contains too few telomeric reads per gigabase
sequenced for routine clinical use, so we developed a telomere capture
preparation in which a biotinylated oligonucleotide complementary to the
3$'$ overhang, combined with restriction digestion of non-telomeric DNA,
enriches telomeric reads by several thousand-fold. A head-to-head
comparison in HEK293T DNA demonstrated that the length distribution is
conserved between telomere capture and whole-genome libraries, while
telomere capture dramatically increases telomeres recovered per gigabase
sequenced.

Digital mean telomere lengths were tightly correlated with TRF
measurements across 14 matched samples and with flow-FISH in 7 RTEL1
mutant samples. Bootstrapping the standard error of the mean as a
function of the number of measured telomeres produced an exponentially
decaying curve that plateaued at a maximal precision of 30--40 bp,
establishing a concrete resolution limit for the assay. TRF systematically
overestimated the digital mean by 1000--3000 bp across all samples,
most likely because undigested subtelomeric sequence inflates the
restriction fragment; flow-FISH overestimated by $\sim$1500 bp on
average. Digital measurement confidently separated wild-type and
pathologically short PARN knock-out hESCs ($p < 2 \times 10^{-16}$,
Wilcoxon rank-sum test), and resolved the expected shortening in TIN2
T284R heterozygous and homozygous hESCs relative to the wild-type
parental line.

Together these results show that digital telomere measurement delivers
individual-telomere resolution, chromosome anchoring and a well-defined
precision floor that no existing assay achieves. With this foundation in
place, we next asked whether the method could track known telomere
dynamics in cultured cells.

\subsection{Digital telomere measurement captures attrition and de novo elongation in cultured cells}

To validate that the method tracks telomere shortening in the absence
of telomerase, we sequenced CDKN2A$^{-/-}$ TERT$^{-/-}$
AAVSI:hTERT$^{\text{flox}}$ hESCs at 66, 78, 98 and 105 days after
Cre-mediated TERT inactivation. Mean telomere length declined linearly
across passages at a rate of approximately 40 bp per day
(Figure~\ref{fig:tert}), consistent with the expected attrition of
telomerase-null replicating cells and with earlier TRF measurements in
the same genetic background. TIN2 T284R heterozygous and homozygous
hESCs exhibited progressively shorter distributions relative to
wild-type, consistent with a dominant-negative effect on shelterin
\citep{palm2008howshelterin,hockemeyer2015control}.

We next asked whether the method can detect de novo telomere elongation
driven by telomerase over-expression. HEK293T cells were transiently
transfected with plasmids expressing wild-type hTR, the TSQ1 variant
template, or GFP, each paired with TERT. Three days after transfection,
TRAP assays confirmed increased telomerase activity, and telomere
capture sequencing with canonical or TSQ1-targeted oligonucleotides
recovered a $\sim$1000 bp increase in mean telomere length. Variant
capture from cells that had not received TSQ1 did not yield a viable
sequencing library, confirming that variant-specific enrichment is
dependent on the presence of the variant repeat in the sample.

Chromosome-anchored analysis of the transfected cells revealed that
chromosomes whose first-quartile telomere length in the control was
shortest underwent the largest elongation under telomerase
over-expression. This provides single-molecule evidence of preferential
telomerase action at short chromosomes in human cells, extending earlier
observations in yeast and in mammalian systems
\citep{teixeira2004intertelomere,hemann2001shortest,britt2012preferential}.
The detection of de novo elongation by both wild-type and variant
capture further confirms that the pipeline measures full-length, intact
telomeres rather than truncated read fragments.

These culture experiments establish that digital telomere measurement
reproduces the expected attrition and elongation phenotypes with
chromosome-level resolution. Having established the method's quantitative
behaviour in cells, we next applied it to human peripheral blood.

\subsection{High-resolution telomere measurement distinguishes human aging and disease}

\begin{figure}[htbp]
\centering
\includegraphics[width=\textwidth]{figures/fig3_schematic.pdf}
\caption{Schematic of the digital telomere measurement pipeline.
Genomic DNA is either subjected to telomere capture (a biotinylated
oligonucleotide complementary to the 3$'$ overhang combined with
restriction digestion) or processed as whole-genome input. The
sequencing library is run on an Oxford Nanopore device. Reads are
aligned by Telometer to chromosomal termini of the telomere-to-telomere
(T2T) human reference. Per-read telomere lengths are aggregated into
per-chromosome distributions and fed to the classifier.}
\label{fig:schematic}
\end{figure}

\begin{figure}[htbp]
\centering
\includegraphics[width=0.65\textwidth]{figures/fig2_tert_attrition.pdf}
\caption{Telomere shortening in TERT-null hESCs. Mean telomere length
decreases linearly with days post Cre-mediated TERT inactivation at
$\sim$40 bp per day across passages sampled at 66, 78, 98 and 105 days.}
\label{fig:tert}
\end{figure}

To test whether the method can distinguish healthy aging from disease,
we sequenced PBLs from 14 healthy donors aged 18--77. Mean, median,
first-quartile (Q1) and third-quartile (Q3) telomere lengths all
declined with donor age, with the mean and median decreasing by
approximately 24 bp per year (Figure~\ref{fig:aging}). Strikingly, Q3
declined more steeply than Q1, indicating that longer telomeres are lost
more rapidly than shorter ones. Cohort aggregation into young
(18--20, $n{=}5$), middle (35--65, $n{=}6$) and elder ($>$70, $n{=}3$)
groups confirmed a progressive loss of long telomeres and an
accumulation of shorter fractions with age, reinforcing earlier TeSLA
observations in a quantitatively direct way \citep{lai2018comparison}.

\begin{figure}[htbp]
\centering
\includegraphics[width=0.7\textwidth]{figures/fig4_aging_slopes.pdf}
\caption{Age-associated telomere attrition in PBLs from 14 healthy donors
(ages 18--77). Mean and median telomere lengths decline by $\sim$24 bp
per year; the third quartile shortens more steeply than the first
quartile.}
\label{fig:aging}
\end{figure}

We next asked whether the distribution structure distinguishes
telomere biology disorder samples from healthy donors. Seven PBL
samples from 6 individuals aged 6--50 carrying an RTEL1 variant
associated with TBD, together with paired PBL (TB32-B) and bone marrow
(TB32-M) DNA from a Stanford patient harbouring a novel TINF2
p.Gln298fs frameshift, CMML and pulmonary fibrosis, had significantly
shorter telomeres than age-matched healthy donors
\citep{walne2013constitutional,savage2008tinf2}. Digital mean lengths
in the RTEL1 samples correlated strongly with independent flow-FISH
measurements at $R^{2} = 0.84$. Unlike flow-FISH or TRF, the digital
distribution separated the TB32 bone marrow sample, whose telomeres
were systematically longer than those in the matched peripheral blood.
A principal component analysis of the full length distributions
separated the 14 healthy samples from the 9 TBD variant carrier
samples, and the two shortest 1000-bp bins (0--2000 bp) were visibly
enriched in TBD samples.

Motivated by this separation, we asked whether a simple classifier
could detect TBD status from the distribution and age alone. A binary
classifier trained at a 7:3 training-to-test ratio achieved 95\%
accuracy (AUC $=$ 0.98) distinguishing healthy donors from symptomatic
patients, 91\% accuracy (AUC $=$ 0.93) for healthy vs asymptomatic
carriers, and 86\% accuracy (AUC $=$ 0.91) for healthy vs the combined
TBD group (Figure~\ref{fig:classifier}, Table~\ref{tab:classifier}).
Because asymptomatic carriers are often indistinguishable from healthy
individuals by flow-FISH \citep{alder2018diagnostic}, the carrier-
versus-healthy AUC is particularly noteworthy: it suggests that the
full distribution contains signal that mean-based assays discard.

\begin{figure}[htbp]
\centering
\includegraphics[width=0.75\textwidth]{figures/fig1_classifier.pdf}
\caption{Binary classifier performance. Using telomere length
measurements and donor age as features, the classifier achieves 95\%
accuracy (AUC $=$ 0.98) for healthy vs symptomatic patients, 91\%
accuracy (AUC $=$ 0.93) for healthy vs asymptomatic carriers, and 86\%
accuracy (AUC $=$ 0.91) for healthy vs the combined TBD group.}
\label{fig:classifier}
\end{figure}

The interquartile asymmetry observed in healthy aging was abolished in
samples with impaired telomere maintenance. In TERT-null hESCs the
daily rate of shortening was at least double the annual attrition of
the analogous summary statistic in healthy PBLs, and Q1 and Q3
shortened at comparable rates---the opposite of the healthy-aging
pattern. In the RTEL1 cohort the steeper decline of Q3 relative to Q1
was likewise absent. Together these observations suggest that the
interquartile structure of the telomere length distribution encodes
information about the state of the telomere maintenance apparatus, and
that limiting telomerase---rather than negative selection against cells
with very short telomeres---is the dominant driver of the accumulation
of shorter telomeres with age.

\subsection{Telomere length distributions in patient-matched colorectal carcinoma and benign epithelia}

To test whether the method transfers to independent, whole-genome
long-read data, we re-analysed 30X on average whole-genome sequencing
from 20 colorectal carcinoma patients with matched benign and
malignant colonic tissue, predominantly aged 50--70
\citep{barthel2017systematic}. In the benign colonic epithelium, mean,
median, Q1 and Q3 all declined with donor age, again with Q3 shortening
more steeply than Q1, mirroring the pattern we observed in PBLs and
consistent with telomerase-limited attrition in a distinct stem cell
compartment.

Tumour telomere length distributions, in contrast, showed no
correlation with donor age for any summary statistic, as expected for
tumour populations immortalised by reactivated telomerase
\citep{chiba2015cancer,barthel2017systematic}. A per-patient Wilcoxon
rank-sum comparison of benign and malignant distributions showed that
tumour telomeres were significantly shorter than matched benign
epithelia in 75\% of the 20 patients (Figure~\ref{fig:crc}), consistent
with earlier studies of colorectal cancer and TCGA cohorts using TRF,
Q-FISH or TelSeq
\citep{ding2014telseq,demanelis2020determinants,barthel2017systematic,roka2023tumor}.

\begin{figure}[htbp]
\centering
\includegraphics[width=0.6\textwidth]{figures/fig5_crc_tumor_vs_benign.pdf}
\caption{Across 20 patients with matched benign colonic epithelium and
colorectal carcinoma, tumour telomeres are significantly shorter than
benign telomeres in 75\% of patients (Wilcoxon rank-sum test per
patient).}
\label{fig:crc}
\end{figure}

These results extend the clinical reach of the method beyond TBD
diagnostics. The contrast between age-dependent shortening in benign
epithelium and age-independent, short but stable tumour distributions
points to a tunable telomeric rheostat as the dominant regulatory
feature of malignant cells
\citep{palm2008howshelterin,hockemeyer2015control}. Having the full
distribution rather than a single mean is essential to detect and
characterise this phenomenon, because a single mean cannot separate a
short-but-stable rheostat set-point from a rapidly attriting
distribution with a similar mean. Moreover, the observation that
telomerase activity in colonic epithelial stem cells---like that in
hematopoietic stem cells---is insufficient to indefinitely sustain the
longest telomeres argues that age-associated Q3 erosion is a general
feature of telomerase-limited human tissues rather than a PBL-specific
phenomenon.

\subsection{Tables}

\begin{table}[htbp]
\centering
\caption{Qualitative comparison of telomere measurement methods. Only
numeric values that appear in the inputs are listed as numbers; other
attributes use categorical labels.}
\label{tab:methods}
\small
\begin{tabular}{lccl}
\toprule
Method & Single-molecule? & Mean overestimate vs digital (bp) & Main limitation \\
\midrule
TRF (Southern blot)        & no  & 1000--3000 & coarse mean; subtelomere bias \\
Flow-FISH                  & no  & $\sim$1500 & relative signal; PBLs mainly \\
qPCR                       & no  & N/A        & relative signal only \\
STELA                      & yes & N/A        & PCR-limited; short ends \\
TeSLA                      & yes & N/A        & PCR-limited; short ends \\
TelSeq (short-read WGS)    & no  & N/A        & relative content only \\
Digital nanopore (this work) & yes & 0 (reference) & requires long-read sequencer \\
\bottomrule
\end{tabular}
\end{table}

\begin{table}[htbp]
\centering
\caption{Binary classifier performance for three healthy-vs-TBD comparisons,
trained on telomere measurements and donor age at a 7:3 training-to-test ratio.}
\label{tab:classifier}
\begin{tabular}{lccc}
\toprule
Comparison & Accuracy (\%) & AUC & Training:test split \\
\midrule
Healthy vs symptomatic patients     & 95 & 0.98 & 7:3 \\
Healthy vs asymptomatic carriers    & 91 & 0.93 & 7:3 \\
Healthy vs combined TBD group       & 86 & 0.91 & 7:3 \\
\bottomrule
\end{tabular}
\end{table}

\section{Discussion}

For decades, human telomere measurement has returned either a relative
telomeric content or an estimate of the population mean length.
Long-read sequencing is now reshaping access to the chromosomal
terminus, and our digital telomere measurement pipeline exploits this
opportunity to deliver high-resolution, high-throughput measurements of
individual intact telomeres. In this work, the method recapitulated the
dramatic shortening expected in three TBD genotypes both in culture and
in vivo, captured de novo telomere addition under forced telomerase
expression, and detected progressive attrition in human aging with and
without functional telomerase. Along the way, we report a previously
undescribed TINF2 p.Gln298fs frameshift in a Stanford patient with
CMML, pulmonary fibrosis and short-for-age telomeres, directly
illustrating how distribution-level readouts can deepen a TBD workup
\citep{savage2008tinf2,alder2018diagnostic}.

Compared to TRF and flow-FISH, digital telomere measurement is
information-rich, requires less input DNA and avoids the polymerase-
processivity limits of STELA and TeSLA
\citep{baird2003extensive,lai2018comparison}. Our data show that TRF
systematically overestimates mean telomere length by up to several
thousand base pairs, a bias that is best explained by incomplete
subtelomeric digestion. Flow-FISH is closer to the digital mean but
still overestimates by $\sim$1500 bp on average, and its utility is
largely restricted to PBLs. Chromosome-level anchoring against the T2T
reference further enables analyses that no prior method supports: we
show that chromosomes with the shortest first-quartile telomeres at
baseline undergo the largest elongation under telomerase over-expression,
providing single-molecule evidence for the preferential action of
telomerase on the shortest telomeres in human cells
\citep{teixeira2004intertelomere,hemann2001shortest,britt2012preferential}.

A second biological contribution of this work is the use of the full
length distribution rather than only its mean. In healthy PBLs and in
benign colonic epithelium, the third quartile shortens faster than the
first---a pattern consistent with insufficient telomerase in the stem
cell compartment to sustain the longest telomeres rather than with
selection against cells carrying very short telomeres
\citep{britt2012preferential}. In TERT-null hESCs and in RTEL1 carriers,
this asymmetry collapses, because telomere maintenance is already
globally limiting. Interquartile structure therefore carries
mechanistic information about the telomere maintenance apparatus that
would be invisible to any mean-based assay, and its readout should
generalise to other tissues and to longitudinal designs.

The clinical implications are concrete. Current TBD diagnostics relies
on flow-FISH, age-adjusted norms and targeted exon sequencing
\citep{alder2018diagnostic,armanios2012telomeres}. Symptomatic patients
often present with mean telomere lengths at or below the first
percentile for age, but TBDs also display genetic anticipation, so
asymptomatic carriers of the same mutation may fall in the lower
quintile rather than the first percentile and go undiagnosed for
decades \citep{walne2013constitutional,savage2008tinf2}. Our classifier
distinguished healthy donors from asymptomatic RTEL1 carriers with
AUC $=$ 0.93 using only the telomere length distribution and age as
features. Because asymptomatic carriers are notoriously difficult to
distinguish from healthy individuals by flow-FISH, this result fills a
real diagnostic gap and suggests that digital telomere measurement is
suitable to complement, rather than simply replace, the current triage.

A third contribution is the perspective this method provides on the
cancer telomere rheostat. In 20 colorectal carcinoma patients, tumour
telomere length distributions were shorter than matched benign
epithelia in 75\% of cases and displayed no correlation with donor age,
consistent with telomerase-stabilised set-points in the tumour
\citep{chiba2015cancer,barthel2017systematic,roka2023tumor}. Because
tumour telomere distributions are known to be stable over time, the
set-point must reflect a dynamic equilibrium between division rate,
telomerase activity and shelterin dosage that has, until now, been
difficult to read out at single-molecule resolution. Digital telomere
measurement offers a direct window on this equilibrium and is a natural
tool for systematic perturbation studies of the rheostat.

A surprising observation across every experiment we performed is that
measurable telomeres as short as tens of base pairs are present in
every sample---including wild-type hESCs. That the shortest fractions
are nonetheless sensitive to telomere maintenance perturbations
suggests that cells can tolerate some quantity of very short telomeres
without triggering replicative senescence or chromosomal instability,
likely through shelterin-mediated protection
\citep{palm2008howshelterin}. Future work using this pipeline could
quantify the critical mass of very short telomeres or the shelterin
threshold at which senescence or telomere fusions are triggered.

This study has several limitations. The cohorts analysed here are
cross-sectional, with a modest TBD sample size ($n{=}9$ variant
carriers) and a colorectal cohort of 20 patients; longitudinal data
would be required to demonstrate the prognostic value of individual
patient trajectories. Capture efficiency may vary with telomere state,
and variant-specific capture is obligately dependent on sequence
complementarity. Finally, the classifier is a simple binary model
trained on a small number of samples and should be validated in larger,
more diverse cohorts before it enters clinical decision-making.

Looking ahead, digital telomere measurement is poised to become a
general-purpose biomarker platform for aging and disease. As long-read
sequencing moves into routine clinical genomics
\citep{nurk2022complete}, analyses of telomere length distributions
from larger, more diverse, and ideally longitudinal studies will let
telomere length become a robust and perhaps even predictive biomarker
of aging, TBD progression, and cancer biology. Integration with
tissue-resolved aging panels, single-cell readouts and quantitative
models of the telomere rheostat will transform our ability to decode
human aging from a single chromosome-anchored measurement.

\bibliographystyle{unsrtnat}
\bibliography{references}

\end{document}
